% Template for post or presentation abstracts submitted to ILAS 2022
% 1. Fill in the presenter's information (name, affiliation, email
% address) and talk information (title of presentation, abstract).
% 2. Compile to make sure there are no errors.
% 3. Save the .tex file for your abstract with a distinctive name,
% ideally "FamilynameGivenname.tex".
% 4. Upload your .tex file to https://clr.ie/129557
%       
% * Please keep your LaTeX commands simple, and avoid the use of
%    user-defined macros.
% * Include details of co-authors and funding acknowledgements after the abstract.
% * Upload only the LaTeX file (with .tex extension).
% * Questions: Contact Niall Madden (Niall.Madden@NUIGalway.ie)

\documentclass[a4paper]{amsart} 
\usepackage{amssymb} 
\usepackage{hyperref}
\pagestyle{empty}
\date{} 

%% IGNORE THE NEXT 4 LINES
\makeatletter % 
\def\labelenumi{\theenumi.}
\def\theenumi{\@arabic\c@enumi}
\makeatother
%% STOP IGNORING NOW

\begin{document}

\title{Majorization and \\ Triangular Polynomial Matrices}
\author{Richard Hollister} %Presenting author only
\address{University at Buffalo} % Affiliation of presenting author
\email{rahollis@buffalo.edu} % Email address of presenting author only

\maketitle

% The abstract  ***no more than one page***

Majorization is a partial ordering of $ \mathbb{R}^n $ with numerous applications across many areas of study, \cite{Marshall-Olkin-I:TMA,Arnold-MHTE}.  This classical concept can trace its origins back to an equivalent ordering given by Muirhead in 1903, \cite{Muirhead-SAFA}.  In this talk, we discuss recent work highlighting the connection between majorization and the diagonals of triangular matrix polynomials.  The question being addressed was first considered in the PhD thesis of Eduardo Marques de S\'a, \cite{MarquesdeSa-IMEFI}:  which polynomial diagonals are possible in a triangular realization of a given Smith form?  The current work answers the same question, but in a conceptually and computationally simpler way using majorization.  The end result is an implementable algorithm that can be used to compute such a triangular realization.

This work relates to research investigating the triangularization of matrix polynomials.  It was shown in \cite{TTZ-TMP} that every regular matrix polynomial over an algebaically closed field can be triangularized.  In more recent work by Anguas, Dopico, Hollister, and Mackey, it was shown that every regular matrix over an \emph{arbitrary field} can be quasi-triangularized with diagonal blocks having bounded sizes, \cite{ADHM-QRMP}.  The role that majorization plays in these results is also discussed.

%% Coauthor details here (optional - delete if not applicable)
% and funding acknowledgment (optional - delete if not applicable)
\emph{This is joint work with Luis Miguel Anguas (Universidad Polit\'ecnica de Madrid), Froil\'{a}n Dopico (Universidad Carlos III de Madrid), and D. Steven Mackey (Western Michigan University).
Supported by ``Proyecto de I+D+i PID2019-106362GB-I00 financiado por MCIN/AEI/10.13039/501100011033'' and ``Ministerio de Econom\'ia, Industria y Competitividad (MINECO)'' of Spain through grants MTM-2015-65798-P and BES-2013-065688.}

\begin{thebibliography}{1}

\bibitem{ADHM-QRMP}
L. M. Anguas, F. Dopico, R. Hollister, D. S. Mackey.
\newblock Quasi-triangularization of regular matrix polynomials over arbitrary fields.
\newblock In review, 2021.

\bibitem{Arnold-MHTE}
C. R. Arnold.
\newblock Majorization: Here, There, and Everywhere.
\newblock {\em Statistical Science}, 22(3):407-413, 2007.

\bibitem{MarquesdeSa-IMEFI} 
E. Marques de S\'a.
\newblock Imers\~ao de Matrizes e Entrela\c{c}amento de Factores Invariantes.
\newblock {\em Doctoral Dissertation}, Universidade de Coimbra, 1979.

\bibitem{Marshall-Olkin-I:TMA} 
A. W. Marshall, I. Olkin, B. C. Arnold.
\newblock Inequalities: Theory of Majorization and Its Applications.
\newblock {\em Springer}, 2011.

\bibitem{Muirhead-SAFA}
R. F. Muirhead.
\newblock Some methods applicable to identities and inequalities of symmetric algebraic functions of $n$ letters.
\newblock {\em Proc. Edinburgh Math. Soc.}, 21:144-157, 1903.

\bibitem{TTZ-TMP}
L. Taslaman, F. Tisseur, and I. Zaballa.
\newblock Triangularizing matrix polynomials.
\newblock {\em Linear Algebra Appl.}, 439:1679-1699, 2013.



\end{thebibliography}


\end{document}


% Please leave the next bit alone  
%%% Local Variables: 
%%% mode: latex
%%% TeX-master: "../ILAS-2022"
%%% End:
